\chapter{KAJIAN PUSTAKA}

\section{\textit{Invoice Tracking}}

\textit{Invoice tracking} adalah pemantauan dokumen \textit{invoice} berdasarkan identitas, status, waktu, dan riwayat perpindahannya. Dokumen diperlakukan sebagai objek proses yang berubah dari satu tahapan ke tahapan berikutnya. Pencatatan tersebut membentuk \textit{audit trail} yang menunjukkan posisi dokumen, waktu perubahan status, dan tindakan yang telah dilakukan. Teknologi pelacakan dokumen mendukung transparansi aliran dokumen dan pengendalian proses administrasi \cite{10.1007/978-3-642-55032-4_33}.

Dalam konteks penagihan, informasi pelacakan membantu pengguna mengetahui \textit{invoice} yang telah diterima, sedang diproses, atau memerlukan tindak lanjut. Analisis data \textit{invoice} juga dapat digunakan untuk mendukung prioritisasi sumber daya administrasi \cite{SchoonbeeEtAl2022}. Literatur klasik menempatkan sistem pendukung keputusan sebagai sarana yang membantu proses keputusan \cite{ShimEtAl2002}. Literatur kontemporer memperluas posisi tersebut dengan menekankan pembagian peran manusia dan kecerdasan artifisial dalam sistem sosio-teknis \cite{StoreyEtAl2024}. Oleh karena itu, rekomendasi prioritas dalam penelitian ini menjadi informasi pendukung bagi admin, bukan keputusan otomatis yang menggantikan kewenangan pengguna.

\section{\textit{Proof of Delivery}}

\textit{Proof of Delivery} (POD) merupakan bukti bahwa barang atau dokumen telah diterima oleh pihak tujuan. POD dapat berbentuk dokumen fisik atau rekaman digital yang memuat informasi penerima, waktu penerimaan, dan konfirmasi serah terima. Dalam proses logistik, keandalan POD diperlukan untuk memperjelas penyelesaian pengiriman dan mengurangi ketidakpastian status layanan \cite{gunadhya_rajawali_logistik_2021}.

POD melengkapi informasi pelacakan karena status pengiriman menunjukkan tahapan proses, sedangkan POD mengonfirmasi bahwa penerimaan telah terjadi. Perkembangan POD digital mencakup pemindaian dokumen, dukungan pengenalan dokumen, dan integrasi dengan infrastruktur digital lain \cite{LinEtAl2026POD}. Dalam penelitian ini, POD menjadi bagian dari riwayat pengiriman \textit{invoice} dan konteks penerapan hasil klasifikasi, bukan atribut yang membentuk \textit{ground truth}.

\section{\textit{Rule-Based System}}

\textit{Rule-Based System} merepresentasikan pengetahuan dalam sekumpulan aturan yang umumnya berbentuk IF--THEN. Sistem ini terdiri atas fakta masukan, basis aturan, dan mekanisme inferensi yang mencocokkan fakta dengan kondisi aturan untuk menghasilkan keluaran \cite{BuchananShortliffe1984,LiuEtAl2016RuleBased}. Aturan dapat berasal dari pengetahuan ahli, dokumen kebijakan, atau prosedur operasional yang telah dinyatakan secara eksplisit.

Keterlacakan menjadi karakteristik penting ketika setiap keluaran dapat dihubungkan dengan aturan yang aktif. Sifat tersebut mendukung pemeriksaan alasan keputusan dan kesesuaian dengan kebijakan operasional. Namun, basis aturan yang besar dapat mengandung redundansi dan menjadi sulit dipelihara, sehingga efektivitas setiap aturan perlu diperiksa \cite{BenDavid2008}. Literatur interpretabilitas menegaskan bahwa transparansi harus dibuktikan melalui representasi dan jejak keputusan yang benar-benar dapat diperiksa, bukan diasumsikan hanya karena model menggunakan aturan \cite{Rudin2019,RudinEtAl2022}.

\section{\textit{Knowledge Acquisition}}

\textit{Knowledge Acquisition} adalah proses memperoleh pengetahuan domain dari sumber yang relevan untuk membangun sistem berbasis pengetahuan. Sumber tersebut dapat berupa ahli, SOP, dokumen kebijakan, observasi proses, dan catatan historis. Rekayasa pengetahuan memandang proses ini sebagai kegiatan pemodelan pengetahuan domain, bukan sekadar memindahkan informasi ke dalam program \cite{StuderEtAl1998}.

Hasil akuisisi perlu dinyatakan sebagai konsep, kondisi, batas, pengecualian, dan hubungan yang dapat diperiksa. Pada sistem berbasis aturan, hasil tersebut menjadi bahan untuk menyusun basis aturan dan mekanisme penjelasan \cite{BuchananShortliffe1984}. Dalam penelitian ini, \textit{Knowledge Acquisition} digunakan untuk memperoleh ketentuan SOP yang diperlukan oleh \textit{Rule-Based}. Proses tersebut tidak digunakan untuk membentuk label admin atau mengubah \texttt{expert\_label}.

\section{\textit{Rule Construction}}

\textit{Rule Construction} mengubah hasil akuisisi pengetahuan menjadi aturan yang memiliki kondisi, keluaran, dan urutan penerapan yang jelas. Konstruksi aturan mencakup penetapan atribut yang diperiksa, operator logika, ambang, aturan \textit{default}, serta penanganan ketika beberapa aturan aktif. Setiap aturan perlu memiliki makna operasional dan sumber yang dapat ditelusuri agar hubungan antara SOP dan basis aturan tetap dapat diperiksa \cite{StuderEtAl1998}.

Pemeriksaan aturan mencakup konsistensi keluaran, ketercakupan kondisi, konflik, tumpang tindih, dan redundansi. Verifikasi menilai apakah aturan dibangun sesuai spesifikasi, sedangkan validasi menilai apakah aturan memadai untuk tujuan domain \cite{Meseguer1992}. Redundansi juga perlu dinilai dari kontribusi aturan terhadap perilaku keseluruhan sistem dan kemudahan pemeliharaan \cite{BenDavid2008}. Basis aturan harus dibekukan sebelum evaluasi terhadap data uji agar label uji tidak memengaruhi konstruksi aturan.

\section{\textit{Decision Tree}}

\textit{Decision Tree} adalah model klasifikasi yang membagi ruang fitur secara rekursif. Setiap simpul internal memuat pengujian terhadap fitur, setiap cabang menunjukkan hasil pengujian, dan setiap daun memberikan kelas prediksi. Pemisahan dipilih untuk membentuk kelompok yang semakin homogen terhadap kelas target \cite{Quinlan1986,JamesEtAl2023}. Kedalaman pohon dan ukuran minimum daun dapat dibatasi untuk mengendalikan kompleksitas dan risiko \textit{overfitting}. Walaupun pohon termasuk model yang secara intrinsik dapat diperiksa, keterperiksaan praktis tetap bergantung pada kompleksitas struktur dan kejelasan jalur keputusannya \cite{RudinEtAl2022}.

CART membentuk pohon melalui pemisahan biner dan dapat digunakan untuk klasifikasi maupun regresi \cite{BreimanEtAl1984}. Scikit-learn menyediakan implementasi pohon keputusan dengan pendekatan CART yang dioptimalkan dan pemisahan biner \cite{PedregosaEtAl2011,ScikitLearnTreeGuide}. Penelitian ini menggunakan \texttt{DecisionTreeClassifier} dengan \texttt{criterion="entropy"}. Kriteria tersebut menentukan ukuran ketidakmurnian, tetapi tidak mengubah implementasi menjadi C4.5. Oleh karena itu, istilah yang digunakan adalah \textit{Decision Tree} berbasis \textit{entropy}, bukan algoritma C4.5.

\section{\textit{Entropy}}

\textit{Entropy} mengukur ketidakmurnian distribusi kelas pada suatu himpunan data. Untuk himpunan $S$ dengan $c$ kelas, \textit{entropy} dirumuskan sebagai berikut \cite{Quinlan1986,JamesEtAl2023}.

\begin{equation}
H(S)=-\sum_{i=1}^{c}p_i\log_2 p_i,
\label{eq:entropy-bab-dua}
\end{equation}

dengan $p_i$ sebagai proporsi observasi pada kelas ke-$i$. Nilai \textit{entropy} sama dengan nol ketika seluruh observasi berada pada satu kelas. Nilainya meningkat ketika distribusi kelas menjadi lebih beragam.

Kualitas kandidat pemisahan dapat dilihat dari penurunan ketidakmurnian setelah data dibagi ke dalam $m$ simpul anak. Penurunan \textit{entropy} dirumuskan sebagai berikut.

\begin{equation}
\Delta H=H(S)-\sum_{j=1}^{m}\frac{|S_j|}{|S|}H(S_j),
\label{eq:penurunan-entropy}
\end{equation}

dengan $S_j$ sebagai himpunan observasi pada simpul anak ke-$j$. Kandidat dengan penurunan ketidakmurnian lebih besar lebih disukai, dengan tetap memperhatikan batas kompleksitas model. Pada scikit-learn, \textit{entropy} merupakan salah satu kriteria pemisahan dalam implementasi CART biner \cite{ScikitLearnTreeGuide}.

\section{Perbedaan ID3, C4.5, dan CART}

ID3, C4.5, dan CART merupakan keluarga algoritma pohon keputusan yang memiliki prosedur berbeda. Perbedaannya perlu dijelaskan agar kriteria \textit{entropy} tidak disamakan dengan implementasi C4.5. Ringkasan perbedaannya disajikan pada Tabel \ref{tab:perbandingan-algoritma-pohon}.

\begin{table}[htbp]
    \centering
    \caption{Perbedaan ID3, C4.5, CART, dan Implementasi Penelitian}
    \label{tab:perbandingan-algoritma-pohon}
    \small
    \renewcommand{\arraystretch}{1.2}
    \begin{tabularx}{\textwidth}{p{0.15\textwidth}p{0.18\textwidth}p{0.20\textwidth}X}
        \toprule
        \textbf{Metode} & \textbf{Bentuk Pemisahan} & \textbf{Kriteria Utama} & \textbf{Karakteristik} \\
        \midrule
        ID3 & Umumnya multi-cabang pada fitur kategorikal & \textit{Information gain} berbasis \textit{entropy} & Algoritma induksi pohon awal; tidak menjadi implementasi penelitian. \\
        C4.5 & Mendukung fitur kategorikal dan numerik & \textit{Gain ratio} & Pengembangan ID3 dengan penanganan nilai kontinu, data hilang, dan pemangkasan. \\
        CART & Pemisahan biner & Penurunan impurity; lazimnya Gini untuk klasifikasi & Mendukung klasifikasi dan regresi serta pemangkasan berbasis kompleksitas. \\
        Implementasi penelitian & Pemisahan biner scikit-learn & \texttt{criterion="entropy"} & Menggunakan CART yang dioptimalkan dengan kriteria \textit{entropy}; bukan C4.5. \\
        \bottomrule
    \end{tabularx}
\end{table}

ID3 dijelaskan oleh Quinlan melalui induksi pohon dan \textit{information gain} \cite{Quinlan1986}. C4.5 mengembangkan pendekatan tersebut dengan \textit{gain ratio} dan prosedur tambahan \cite{Quinlan1993C45}. CART dikembangkan sebagai kerangka pohon klasifikasi dan regresi dengan pemisahan biner \cite{BreimanEtAl1984}. Dokumentasi scikit-learn menegaskan bahwa implementasinya menggunakan versi CART yang dioptimalkan \cite{ScikitLearnTreeGuide}.

\FloatBarrier

\section{\textit{Supervised Classification}}

\textit{Supervised classification} mempelajari fungsi pemetaan dari fitur ke kelas menggunakan observasi yang telah berlabel. Data latih digunakan untuk membentuk model, sedangkan data yang tidak digunakan dalam pelatihan diperlukan untuk memperkirakan kemampuan prediksi pada kasus baru \cite{JamesEtAl2023}. Dalam klasifikasi biner penelitian ini, label target terdiri atas \textit{Urgent} dan \textit{Not Urgent}.

Pemisahan data harus diterapkan pada seluruh proses yang belajar dari data. Transformasi, pemilihan parameter, dan keputusan pemodelan tidak boleh menggunakan informasi data uji. Penggunaan informasi tersebut menimbulkan \textit{data leakage} dan dapat menghasilkan estimasi performa yang terlalu optimistis \cite{KapoorNarayanan2023}. Label historis admin pada \texttt{expert\_label} menjadi \textit{ground truth} bersama, sedangkan \textit{Rule-Based} dan \textit{Decision Tree} menghasilkan prediksi yang dibandingkan terhadap label tersebut.

\section{\textit{Confusion Matrix}}

\textit{Confusion matrix} membandingkan kelas prediksi dengan \textit{ground truth}. Pada klasifikasi biner, matriks ini terdiri atas \textit{true positive} (TP), \textit{true negative} (TN), \textit{false positive} (FP), dan \textit{false negative} (FN). Definisi tersebut berasal dari landasan evaluasi klasifikasi yang telah mapan \cite{SokolovaLapalme2009} dan tetap menjadi dasar panduan evaluasi \textit{machine learning} kontemporer \cite{RainioEtAl2024}. Struktur matriks yang digunakan ditunjukkan pada Tabel \ref{tab:confusion-matrix-biner}.

\begin{table}[htbp]
    \centering
    \caption{\textit{Confusion Matrix} untuk Klasifikasi Prioritas Invoice}
    \label{tab:confusion-matrix-biner}
    \renewcommand{\arraystretch}{1.2}
    \begin{tabular}{lcc}
        \toprule
        & \multicolumn{2}{c}{\textbf{Prediksi}} \\
        \cmidrule(lr){2-3}
        \textbf{Ground Truth} & \textbf{Urgent} & \textbf{Not Urgent} \\
        \midrule
        \textit{Urgent} & TP & FN \\
        \textit{Not Urgent} & FP & TN \\
        \bottomrule
    \end{tabular}
\end{table}

Kelas \textit{Urgent} ditetapkan sebagai kelas positif. TP adalah \textit{invoice Urgent} yang diprediksi \textit{Urgent}; FN adalah \textit{invoice Urgent} yang diprediksi \textit{Not Urgent}. FP adalah \textit{invoice Not Urgent} yang diprediksi \textit{Urgent}; TN adalah \textit{invoice Not Urgent} yang diprediksi \textit{Not Urgent}. Keempat komponen tersebut menjadi dasar perhitungan metrik dan analisis pola kesalahan.

\FloatBarrier

\section{\textit{Accuracy}}

\textit{Accuracy} mengukur proporsi seluruh prediksi yang benar terhadap jumlah observasi. Rumusnya adalah sebagai berikut \cite{SokolovaLapalme2009,RainioEtAl2024}.

\begin{equation}
\mathrm{Accuracy}=\frac{TP+TN}{TP+TN+FP+FN}.
\label{eq:accuracy-bab-dua}
\end{equation}

Metrik ini memberikan ringkasan performa umum, tetapi tidak menunjukkan jenis kesalahan. Pada distribusi kelas yang tidak seimbang, nilai yang tinggi dapat didominasi oleh kelas mayoritas. Oleh karena itu, \textit{accuracy} perlu dibaca bersama \textit{precision}, \textit{recall}, F1, serta jumlah FP dan FN \cite{RainioEtAl2024}.

\section{\textit{Precision}}

\textit{Precision} mengukur proporsi prediksi \textit{Urgent} yang benar-benar memiliki \textit{ground truth Urgent}. Rumusnya adalah sebagai berikut \cite{SokolovaLapalme2009,RainioEtAl2024}.

\begin{equation}
\mathrm{Precision}_{Urgent}=\frac{TP}{TP+FP}.
\label{eq:precision-bab-dua}
\end{equation}

Nilai \textit{precision} menurun ketika jumlah FP meningkat. Dalam konteks operasional, FP menunjukkan \textit{invoice Not Urgent} yang memperoleh rekomendasi \textit{Urgent}. Metrik ini menjelaskan ketepatan pemberian rekomendasi positif, bukan kemampuan menemukan seluruh kasus \textit{Urgent}.

\section{\textit{Recall}}

\textit{Recall} mengukur proporsi \textit{invoice} dengan \textit{ground truth Urgent} yang berhasil diprediksi \textit{Urgent}. Rumusnya adalah sebagai berikut \cite{SokolovaLapalme2009,RainioEtAl2024}.

\begin{equation}
\mathrm{Recall}_{Urgent}=\frac{TP}{TP+FN}.
\label{eq:recall-bab-dua}
\end{equation}

Nilai \textit{recall} menurun ketika FN meningkat. FN merupakan kasus \textit{Urgent} yang diprediksi \textit{Not Urgent}, sehingga metrik ini digunakan untuk menilai kemampuan metode menemukan kasus yang memerlukan prioritas. Interpretasinya perlu dipisahkan dari \textit{precision} karena kedua metrik menilai pola kesalahan yang berbeda.

\section{\textit{F1-score}}

\textit{F1-score} merupakan rata-rata harmonik \textit{precision} dan \textit{recall}. Rumus F1 untuk kelas \textit{Urgent} adalah sebagai berikut \cite{SokolovaLapalme2009,RainioEtAl2024}.

\begin{equation}
\mathrm{F1}_{Urgent}=2\times\frac{\mathrm{Precision}_{Urgent}\times\mathrm{Recall}_{Urgent}}{\mathrm{Precision}_{Urgent}+\mathrm{Recall}_{Urgent}}.
\label{eq:f1-bab-dua}
\end{equation}

F1 kelas \textit{Urgent} berfokus pada keseimbangan kesalahan untuk kelas positif. Sementara itu, \textit{macro F1} menghitung rata-rata F1 kedua kelas dengan bobot yang sama.

\begin{equation}
\mathrm{Macro\ F1}=\frac{\mathrm{F1}_{Urgent}+\mathrm{F1}_{Not\ Urgent}}{2}.
\label{eq:macro-f1-bab-dua}
\end{equation}

\textit{Macro F1} berguna ketika kedua kelas perlu mendapat perhatian yang setara. Metrik ini tetap tidak menggantikan analisis FP, FN, dan konsekuensi operasional setiap kesalahan.

\section{\textit{Hold-out Validation}}

\textit{Hold-out validation} membagi dataset menjadi data latih dan data uji yang tidak tumpang tindih. Model dibangun menggunakan data latih dan dievaluasi sekali pada data uji \cite{JamesEtAl2023}. Pada klasifikasi dengan distribusi kelas tidak seimbang, stratifikasi digunakan agar proporsi kelas pada kedua bagian tetap mendekati dataset asal.

Kelebihan \textit{hold-out} adalah prosedurnya sederhana dan biaya komputasinya rendah. Keterbatasannya adalah hasil dapat dipengaruhi oleh satu pembagian data, terutama pada dataset kecil \cite{Kohavi1995}. Metode ini menjadi dasar teoritis bagi skenario E1 dan E2, yang menggunakan proporsi data latih dan uji berbeda untuk memeriksa sensitivitas hasil terhadap pembagian data.

\section{\textit{Cross Validation}}

Pada K-fold \textit{cross-validation}, dataset dibagi menjadi $K$ bagian. Setiap bagian digunakan satu kali sebagai data validasi, sedangkan bagian lain digunakan sebagai data latih. Prediksi dari seluruh fold dapat digabungkan menjadi prediksi \textit{out-of-fold} sehingga setiap observasi dinilai ketika tidak menjadi bagian dari data latih \cite{ArlotCelisse2010}. Kajian mutakhir menjelaskan bahwa \textit{cross-validation} mengestimasi rata-rata galat prediksi lintas himpunan pelatihan dan bahwa hasil antar-fold tidak bersifat independen \cite{BatesEtAl2024CV}. Karena itu, variasi antar-fold tidak boleh langsung diperlakukan sebagai replikasi independen atau estimasi ketidakpastian konvensional.

\textit{Stratified cross-validation} mempertahankan proporsi kelas pada setiap fold dan dapat mengurangi distorsi akibat pembagian kelas yang tidak seimbang \cite{Kohavi1995}. Pemilihan parameter harus dilakukan hanya dari data latih pada fold terkait. Prinsip ini mendukung E3, sedangkan hasilnya tetap harus dipahami sebagai estimasi dari data yang tersedia, bukan bukti generalisasi lintas perusahaan.

\section{\textit{Leave-One-Out Cross Validation} (LOOCV)}

LOOCV merupakan bentuk \textit{cross-validation} dengan jumlah fold sama dengan jumlah observasi. Pada setiap iterasi, satu observasi digunakan sebagai data uji dan seluruh observasi lainnya menjadi data latih. Proses diulang hingga setiap observasi memperoleh satu prediksi \textit{out-of-sample} \cite{ArlotCelisse2010}.

LOOCV menggunakan hampir seluruh data untuk pelatihan pada setiap iterasi, tetapi memerlukan pelatihan berulang dan dapat memiliki variasi estimasi yang tinggi \cite{Kohavi1995}. Prosedur ini mendukung E4 sebagai skenario evaluasi tambahan pada dataset kecil. LOOCV tidak otomatis menjadi prosedur terbaik; hasilnya perlu dibaca bersama skenario validasi lain.

\section{\textit{McNemar Test}}

Uji McNemar membandingkan dua keluaran berpasangan pada kasus yang sama \cite{McNemar1947}. Penggunaannya dalam perbandingan \textit{classifier} telah dibahas dalam landasan evaluasi klasik \cite{Dietterich1998}, sedangkan panduan kontemporer tetap menempatkannya sebagai uji untuk dua prediksi berpasangan pada unit observasi yang sama \cite{RainioEtAl2024}. Setiap kasus dinyatakan benar atau salah untuk kedua metode. Susunan kasus berpasangan ditunjukkan pada Tabel \ref{tab:mcnemar-berpasangan}.

\begin{table}[htbp]
    \centering
    \caption{Tabel Kontingensi untuk Uji McNemar Berpasangan}
    \label{tab:mcnemar-berpasangan}
    \renewcommand{\arraystretch}{1.2}
    \begin{tabular}{lcc}
        \toprule
        & \multicolumn{2}{c}{\textbf{Decision Tree}} \\
        \cmidrule(lr){2-3}
        \textbf{Rule-Based} & \textbf{Benar} & \textbf{Salah} \\
        \midrule
        Benar & $n_{11}$ & $b$ \\
        Salah & $c$ & $n_{00}$ \\
        \bottomrule
    \end{tabular}
\end{table}

Nilai $b$ menyatakan kasus yang hanya benar pada \textit{Rule-Based}, sedangkan $c$ menyatakan kasus yang hanya benar pada \textit{Decision Tree}. Kasus $n_{11}$ dan $n_{00}$ bersifat \textit{concordant}; uji ditentukan oleh $b$ dan $c$ yang bersifat \textit{discordant}. Hipotesis nol menyatakan bahwa kedua metode memiliki proporsi kesalahan marginal yang sama.

Untuk jumlah kasus \textit{discordant} yang kecil, \textit{exact McNemar} dua sisi dapat dihitung berdasarkan distribusi binomial sebagai berikut.

\begin{equation}
p_{exact}=\min\left(1,2\sum_{k=0}^{\min(b,c)}\binom{b+c}{k}(0.5)^{b+c}\right).
\label{eq:exact-mcnemar}
\end{equation}

Nilai $p$ di atas taraf signifikansi menunjukkan bahwa bukti belum cukup untuk menolak hipotesis nol. Interpretasi tersebut tidak membuktikan bahwa kedua metode identik dan tidak menghapus kebutuhan untuk membahas perbedaan deskriptif serta pola kesalahan.

\FloatBarrier

% \section{Penelitian Terdahulu}

% Penelitian terdahulu dipetakan berdasarkan domain, metode, target, sumber acuan, prosedur validasi, serta relevansinya terhadap penelitian ini. Pemetaan tersebut disajikan pada Tabel \ref{tab:penelitian_terdahulu}.

% \begin{table}[htbp]
%     \centering
%     \caption{Perbandingan Penelitian Terdahulu}
%     \label{tab:penelitian_terdahulu}
%     \scriptsize
%     \setlength{\tabcolsep}{2pt}
%     \renewcommand{\arraystretch}{1.15}
%     \begin{tabularx}{\textwidth}{p{0.095\textwidth}p{0.075\textwidth}p{0.085\textwidth}p{0.075\textwidth}p{0.09\textwidth}p{0.085\textwidth}p{0.115\textwidth}p{0.115\textwidth}X}
%         \toprule
%         \textbf{Peneliti/Tahun} & \textbf{Domain} & \textbf{Metode} & \textbf{Target} & \textbf{Ground Truth/Data Acuan} & \textbf{Validasi} & \textbf{Temuan Utama} & \textbf{Keterbatasan} & \textbf{Relevansi} \\
%         \midrule
%         Buchanan dan Shortliffe (1984) \cite{BuchananShortliffe1984} & Sistem pakar klinis & \textit{Rule-Based} MYCIN & Rekomendasi klinis & Pengetahuan ahli dan kasus klinis & Evaluasi sistem pakar & Aturan mendukung inferensi dan penjelasan keputusan. & Domain klinis; bukan klasifikasi invoice. & Dasar sistem pakar, akuisisi pengetahuan, dan aturan. \\
%         Sunindyo dkk. (2014) \cite{10.1007/978-3-642-55032-4_33} & Administrasi pemerintah & \textit{Document tracking} & Status dokumen & Rekaman proses dokumen & Evaluasi penerapan teknologi & Pelacakan mendukung transparansi alur administrasi. & Tidak mengevaluasi classifier prioritas. & Dasar konteks \textit{invoice tracking}. \\
%         Ben-David (2008) \cite{BenDavid2008} & \textit{Credit scoring} & Sistem berbasis aturan & Efektivitas aturan & Kasus \textit{credit scoring} nyata & Kontribusi aturan terhadap akurasi & Redundansi dapat dianalisis untuk merampingkan basis aturan. & Satu domain dan satu basis aturan. & Dasar validasi, redundansi, dan pemeliharaan aturan. \\
%         Berhane dkk. (2018) \cite{BerhaneEtAl2018} & Pemetaan lahan basah & DT, RB, dan RF & Kelas lahan basah & Data referensi lapangan & Data \textit{hold-out} dan interval kepercayaan & Membandingkan beberapa classifier pada data uji yang sama. & RB merupakan ruleset turunan model, bukan SOP operasional. & Bukti studi komparatif; domain berbeda. \\
%         Lin dkk. (2026) \cite{LinEtAl2026POD} & Logistik jarak akhir & LCA empat konfigurasi POD & Dampak konfigurasi POD & Data transaksi dan inventori LCA & Perbandingan konfigurasi & POD digital menjadi bagian transformasi proses logistik. & Tidak melakukan klasifikasi prioritas. & Dasar POD digital dan konteks logistik. \\
%         Schoonbee dkk. (2022) \cite{SchoonbeeEtAl2022} & Penagihan invoice & Kandidat metode ML dalam DSS & Interval keterlambatan pembayaran & Data invoice industri & Evaluasi kandidat model & Prediksi invoice digunakan untuk memprioritaskan sumber daya penagihan. & Target pembayaran, bukan pengiriman; tanpa pembanding RB berbasis SOP. & Konteks invoice, ML, dan DSS. \\
%         Firmansyah dan Yulianto (2021) \cite{firmansyah} & Pembayaran invoice & C5.0 & Ketepatan pembayaran & Data historis invoice & Evaluasi klasifikasi & Pohon keputusan diterapkan pada prediksi pembayaran invoice. & Target dan varian algoritma berbeda; tanpa RB. & Studi invoice Indonesia. \\
%         Utami dan Arifyanto (2025) \cite{dyahikbal} & Penagihan listrik & Random Forest & Keterlambatan pembayaran & Data pelanggan historis & K-fold dan \textit{rolling window} & Prediksi mendukung prioritas distribusi tagihan. & Bukan DT versus RB dan bukan prioritas pengiriman. & Konteks prioritas dan validasi multi-skenario. \\
%         Diwandari dkk. (2023) \cite{Diwandari_Sela_Syahputra_Parama_Septiarini_2023} & Bantuan sosial & DT dan AHP & Kelayakan dan prioritas & Data calon penerima & Evaluasi klasifikasi dan pemeringkatan & Klasifikasi dan prioritas dapat menjalankan fungsi berbeda dalam DSS. & Domain berbeda; tidak membandingkan DT dengan RB. & Konteks DSS dan prioritas operasional. \\
%         Shim dkk. (2002) \cite{ShimEtAl2002} & Sistem pendukung keputusan & Kajian teknologi DSS & Dukungan keputusan & Literatur perkembangan DSS & Sintesis konseptual & DSS berkembang sebagai dukungan analitis bagi pengambil keputusan. & Tidak menyediakan eksperimen klasifikasi. & Menetapkan sistem sebagai konteks dukungan keputusan. \\
%         \bottomrule
%     \end{tabularx}
% \end{table}

% Studi pelacakan dokumen dan POD menjelaskan pencatatan status, riwayat proses, dan bukti penerimaan, tetapi tidak mengevaluasi classifier prioritas \cite{10.1007/978-3-642-55032-4_33,LinEtAl2026POD}. Studi sistem pakar dan basis aturan memberikan dasar bagi representasi pengetahuan, validasi, dan pemeliharaan aturan \cite{BuchananShortliffe1984,BenDavid2008}. Keduanya mendukung bagian yang berbeda dari penelitian ini dan tidak dapat diperlakukan sebagai bukti evaluasi klasifikasi invoice.

% Studi komparatif Berhane dkk. membandingkan \textit{Decision Tree}, \textit{Rule-Based}, dan Random Forest pada kasus validasi yang sama, tetapi aturan yang digunakan merupakan ruleset dalam domain citra lahan basah \cite{BerhaneEtAl2018}. Kondisi tersebut berbeda dari \textit{Rule-Based} yang dibangun langsung dari SOP perusahaan. Perbedaan sumber aturan memengaruhi makna transparansi, kepatuhan terhadap SOP, dan prosedur pembaruan.

% Studi invoice yang ditinjau berfokus pada pembayaran atau keterlambatan penagihan \cite{SchoonbeeEtAl2022,firmansyah,dyahikbal}. Studi DSS lain membahas kelayakan dan prioritas pada bantuan sosial \cite{Diwandari_Sela_Syahputra_Parama_Septiarini_2023}. Literatur tersebut menunjukkan relevansi klasifikasi dan prioritisasi, tetapi target, sumber \textit{ground truth}, serta metode pembandingnya tidak sama dengan klasifikasi prioritas pengiriman \textit{invoice}.

% \FloatBarrier

% \section{\textit{Research Gap}}

% Kumpulan penelitian yang ditinjau menunjukkan tiga kelompok pengetahuan. Pertama, pelacakan dokumen dan POD mendukung pemantauan proses serta verifikasi penerimaan. Kedua, sistem berbasis aturan mendukung representasi SOP yang dapat ditelusuri. Ketiga, metode machine learning dan DSS mendukung klasifikasi serta prioritisasi berdasarkan data historis.

% Meskipun demikian, kumpulan penelitian tersebut belum memberikan bukti komparatif langsung mengenai performa, pola kesalahan, dan karakteristik operasional \textit{Rule-Based} berbasis SOP dan \textit{Decision Tree} berbasis data historis admin untuk klasifikasi prioritas pengiriman \textit{invoice}. Studi invoice berfokus pada pembayaran, studi pelacakan dan POD tidak membandingkan classifier, sedangkan studi komparatif \textit{Rule-Based} dan \textit{Decision Tree} menggunakan domain serta sumber aturan yang berbeda.

% Penelitian ini menempatkan label historis admin sebagai \textit{ground truth} bersama. \textit{Rule-Based} dibangun melalui \textit{Knowledge Acquisition}, \textit{Rule Construction}, dan \textit{Rule Validation}, sedangkan \textit{Decision Tree} berbasis \textit{entropy} dilatih dari data historis. Keduanya dibandingkan pada kasus validasi yang sama melalui evaluasi multi-skenario, analisis kesalahan per \textit{invoice}, uji statistik berpasangan, dan analisis karakteristik operasional. Sistem pelacakan \textit{invoice} dan POD tetap ditempatkan sebagai konteks penerapan rekomendasi prioritas.
